\documentclass[12pt,a4paper]{article}
\usepackage[utf8]{inputenc}
\usepackage[T2A]{fontenc}
\usepackage[russian]{babel}
\usepackage{amsmath,amssymb}
\usepackage{array}
\usepackage{longtable}
\usepackage{booktabs}
\usepackage{geometry}
\usepackage{hyperref}
\geometry{margin=2.3cm}

\title{Сравнительный анализ теорий \texttt{TOE.pdf} и \texttt{paper-6.pdf}}
\author{}
\date{}

\begin{document}

\maketitle

\begin{abstract}
В данном документе проводится сравнительный анализ двух теоретических корпусов: спектрально-нелокальной объединяющей программы \texttt{TOE.pdf} и локальной геометрико-спектральной конструкции \texttt{paper-6.pdf}, посвящённой формуле-кандидату для постоянной тонкой структуры. Цель сравнения --- выделить их пересечения, расхождения, различия по уровню строгости, воспроизводимости и научной зрелости, а также дать ориентировочную балльную оценку по единой S2T-совместимой шкале. Баллы ниже не являются ``истиной'', а служат операциональной формой сравнительного суждения на основании изученных текстов.
\end{abstract}

\section{Задача сравнения}

Сравнение строится не по принципу ``какая теория красивее'', а по более строгим вопросам:
\begin{itemize}
    \item что является фундаментальным объектом теории;
    \item каков охват объясняемых явлений;
    \item насколько ясно разведены эвристика и строгий вывод;
    \item насколько теория воспроизводима вычислительно;
    \item насколько она устойчива к проверке в логике S2T;
    \item где теория сильнее как программа, а где сильнее как строго оформленная постановка.
\end{itemize}

\section{Краткая характеристика каждой работы}

\subsection{Корпус \texttt{TOE.pdf}}

\texttt{TOE.pdf} и связанный с ним архив --- это большая объединяющая программа, где фундаментальным объектом выступает нелокальный корреляционный оператор. Из него автор пытается вывести геометрию пространства-времени, гравитацию, калибровочные взаимодействия, массу Хиггса, число поколений, UV-конечность, а далее и широкий феноменологический слой: тёмную материю, бариогенез, модульное время, \(g-2\), реликтовые гравитационные волны и др.

Главное достоинство этой программы --- широта и концептуальная цельность. Главная слабость --- незамкнутость ключевых доказательных узлов, которая признаётся и самим корпусом документов.

\subsection{Работа \texttt{paper-6.pdf}}

\texttt{paper-6.pdf} --- значительно более узкая по замыслу работа. Она рассматривает конкретную формулу-кандидат для \(\alpha^{-1}\) на компактной геометрии \(K = RP^3 \times S^1\) и переводит её в язык KK-редукции, one-loop вычислений, ghost/det$'$, \(\zeta\)-регуляризации и воспроизводимых постановочных проверок.

Главное достоинство \texttt{paper-6.pdf} --- высокий уровень методологической дисциплины: автор ясно разводит эвристику и верификацию, вводит паспорт постановки, описывает, какие части уже получены строго, а какие остаются открытыми. Главная слабость --- узость охвата: работа не является TOE в сильном смысле, а концентрируется на одном физическом числе и ограниченном механизме его геометрико-спектральной мотивации.

\section{Главные пересечения}

Несмотря на разный масштаб, у теорий есть существенные точки пересечения.

\begin{enumerate}
    \item Обе работы исходят из идеи, что наблюдаемая физика может кодироваться в компактной геометрии и спектральных инвариантах.
    \item Обе допускают важную роль ИИ или ИИ-подобной эвристики на этапе генерации кандидатов, но затем пытаются переводить результат в стандартный математико-физический язык.
    \item Обе теории считают, что простое численное совпадение недостаточно и что требуется переход к воспроизводимой вычислительной постановке.
    \item Обе работают на стыке геометрии, спектрального анализа, квантовой теории поля и топологии.
    \item Обе в разной степени признают различие между красивой формулой и строгим выводом.
\end{enumerate}

На этом уровне \texttt{paper-6.pdf} можно рассматривать как частный, более локальный и более дисциплинированный тип программы, тогда как \texttt{TOE.pdf} --- как глобальный и значительно более амбициозный вариант той же общей эпистемологической линии.

\section{Главные расхождения}

\subsection{Масштаб и амбиция}

\texttt{TOE.pdf} претендует на единую теорию широкого класса фундаментальных явлений. \texttt{paper-6.pdf} сосредоточена на одной конкретной величине --- постоянной тонкой структуры --- и на локальной геометрико-спектральной постановке.

Следовательно:
\begin{itemize}
    \item \texttt{TOE.pdf} сильнее по охвату;
    \item \texttt{paper-6.pdf} сильнее по локальной управляемости задачи.
\end{itemize}

\subsection{Статус строгости}

\texttt{paper-6.pdf} гораздо аккуратнее маркирует открытые вопросы. В финале работы прямо перечислены ограничения: не завершённый полный \(\zeta\)-расчёт, незамкнутый механизм линейного \(\pi\)-вклада, модельный статус члена порядка \(S^{-2}_{geo}\), частично проработанная, но незамкнутая стабилизация масштаба \(R\), а также необходимость RG-согласования. Это делает её честной и дисциплинированной с точки зрения S2T.

\texttt{TOE.pdf} и архив тоже содержат самокритику, но в базовых версиях сильнее присутствует язык ``strict derivation'' при не до конца замкнутой доказательной цепочке. Поэтому по формальной дисциплине самопрезентации \texttt{paper-6.pdf} выглядит строже.

\subsection{Воспроизводимость}

\texttt{paper-6.pdf} делает сильный акцент на ``паспорте постановки'', запрете якорения на эксперимент, хранении артефактов и разделении генерации кандидатов от верификации. Это очень близко к духу S2T.

В \texttt{TOE.pdf} воспроизводимость тоже декларируется, но не во всех ключевых местах превращена в явную вычислительную инфраструктуру одинакового уровня детализации.

\subsection{Феноменологический риск}

\texttt{TOE.pdf} выигрывает в том, что ставит на карту много проверяемых следствий. Но эта же широта повышает каскадную зависимость блоков: незавершённость ядра может ослабить множество периферийных выводов.

\texttt{paper-6.pdf} рискует меньше: её локальная задача уже сама по себе трудная, но область утверждений ограничена и лучше контролируется.

\section{Где какая работа строже}

\subsection{Где строже \texttt{paper-6.pdf}}

\begin{itemize}
    \item в разведении эвристической и строгой частей;
    \item в явном описании воспроизводимого протокола;
    \item в фиксации постановочных выборов;
    \item в перечислении незакрытых вычислительных узлов;
    \item в локальной честности статуса результата.
\end{itemize}

\subsection{Где строже или сильнее \texttt{TOE.pdf}}

\begin{itemize}
    \item в масштабе объединяющей концепции;
    \item в широте связи между различными секторами физики;
    \item в числе потенциально фальсифицируемых следствий;
    \item в наличии большой модульной исследовательской программы, а не одного частного кейса.
\end{itemize}

Но здесь важно различать \emph{строгость охвата} и \emph{строгость вывода}. По охвату сильнее \texttt{TOE.pdf}; по локальной дисциплине вывода сильнее \texttt{paper-6.pdf}.

\section{Балльная шкала сравнения}

Вводится ориентировочная шкала от 0 до 10 по каждому критерию:
\begin{itemize}
    \item 0--2: очень слабый уровень;
    \item 3--4: слабый уровень;
    \item 5--6: средний рабочий уровень;
    \item 7--8: сильный уровень;
    \item 9--10: очень сильный уровень.
\end{itemize}

Баллы выражают не ``истинность'' теории, а степень зрелости по конкретному параметру.

\begin{longtable}{p{5.3cm}p{2cm}p{2cm}p{5.6cm}}
\toprule
Критерий & TOE & paper-6 & Комментарий \\
\midrule
\endhead
Масштаб объяснения & 10 & 4 & \texttt{TOE.pdf} претендует на унификацию больших секторов физики; \texttt{paper-6.pdf} решает локальную задачу \(\alpha\). \\
Концептуальная цельность & 8 & 7 & Обе работы внутренне цельные, но TOE сильнее как общая архитектура. \\
Ясность границ утверждений & 6 & 9 & \texttt{paper-6.pdf} гораздо чётче маркирует, что доказано, а что нет. \\
Локальная математическая дисциплина & 6 & 8 & В \texttt{paper-6.pdf} более аккуратная фиксация операторов, детерминантов, \(\zeta\)-регуляризации и постановки. \\
Глобальная доказательная замкнутость & 4 & 5 & Обе работы незамкнуты; у \texttt{paper-6.pdf} задача уже, потому gap контролируемее. \\
Воспроизводимость постановки & 5 & 9 & У \texttt{paper-6.pdf} явно есть паспорт постановки и анти-якорный протокол. \\
Anti-overfit готовность & 4 & 8 & \texttt{paper-6.pdf} значительно ближе к S2T-дисциплине сравнения и запрета подгонки. \\
Честность самокритики & 7 & 9 & Обе работы признают ограничения, но \texttt{paper-6.pdf} делает это строже и компактнее. \\
Фальсифицируемость & 8 & 6 & TOE даёт больше внешних тестов; \texttt{paper-6.pdf} более локальна, но тоже допускает проверку. \\
Модульность программы & 9 & 6 & TOE развёрнута в большой корпус; \texttt{paper-6.pdf} скорее цельная монография по одной задаче. \\
Готовность к S2T-аудиту & 5 & 8 & \texttt{paper-6.pdf} лучше соответствует логике S2T уже на уровне структуры текста. \\
Потенциал научного влияния при успехе & 10 & 7 & Если TOE замкнётся, её влияние будет огромным; у \texttt{paper-6.pdf} влияние уже, но тоже существенно. \\
\midrule
Сумма & 82 & 86 & Сумма отражает зрелость по выбранным критериям, а не ``победу'' по истине. \\
\bottomrule
\end{longtable}

\section{Интерпретация баллов}

На первый взгляд может показаться неожиданным, что при меньшем масштабе \texttt{paper-6.pdf} получает сопоставимую или даже чуть более высокую суммарную оценку. Это не означает, что \texttt{paper-6.pdf} ``глубже'' или ``вернее'' как теория мира. Это означает только, что как \emph{локально дисциплинированная исследовательская постановка} она сейчас выглядит более зрелой.

И наоборот, \texttt{TOE.pdf} как программа значительно мощнее по охвату, но теряет баллы там, где большая амбиция ещё не обеспечена соответствующим уровнем proof-closure, anti-overfit контроля и воспроизводимой формальной фиксации.

\section{Итоговый сравнительный вывод}

Сравнение позволяет сделать следующий аккуратный вывод.

\begin{enumerate}
    \item \texttt{TOE.pdf} сильнее как большая объединяющая программа.
    \item \texttt{paper-6.pdf} сильнее как узкая, честно структурированная и S2T-совместимая исследовательская постановка.
    \item \texttt{TOE.pdf} выигрывает по амбиции, охвату и числу потенциальных внешних тестов.
    \item \texttt{paper-6.pdf} выигрывает по дисциплине формулировок, воспроизводимости, локальной методологической строгости и контролю открытых вопросов.
    \item В терминах научной зрелости ``здесь и сейчас'' \texttt{paper-6.pdf} выглядит более аккуратно доведённой как исследовательский документ.
    \item В терминах потенциальной фундаментальной значимости при условии успешного замыкания сильнее остаётся \texttt{TOE.pdf}.
\end{enumerate}

Иными словами, эти работы не столько конкурируют напрямую, сколько представляют два разных режима теоретической физики. \texttt{TOE.pdf} --- это режим большой унификационной архитектуры. \texttt{paper-6.pdf} --- режим локального, но существенно более дисциплинированного доведения одной гипотезы до проверяемой формы.

\section{Практический вывод для автора}

Если цель состоит в повышении научной зрелости корпуса \texttt{TOE.pdf}, то полезно перенять у \texttt{paper-6.pdf} следующие качества:
\begin{itemize}
    \item паспорт постановки;
    \item жёсткое разведение генерации кандидата и верификации;
    \item явный список незавершённых вычислений;
    \item запрет преждевременной риторики завершённости;
    \item более строгую привязку каждого численного результата к воспроизводимой процедуре.
\end{itemize}

Если же цель состоит в расширении \texttt{paper-6.pdf} до уровня более общей теории, то урок от \texttt{TOE.pdf} противоположный: локальная строгость должна сочетаться с большим архитектурным горизонтом и связью между различными секторами физики.

\section{Заключение}

Сравнение показывает, что \texttt{TOE.pdf} и \texttt{paper-6.pdf} представляют два различных, но взаимно поучительных типа теоретической работы. Первая сильнее по масштабу и объединительной амбиции, вторая --- по локальной строгости и дисциплине исследовательской постановки. Поэтому корректнее не объявлять одну ``лучше'' другой, а признать, что на текущем этапе \texttt{paper-6.pdf} выглядит более зрелой как научный документ, тогда как \texttt{TOE.pdf} остаётся более мощной по потенциальному фундаментальному охвату, но нуждающейся в более строгом доведении.

\end{document}